\section{Лекция 2. Отношения и отображения}
\noeqref{eq:bin_rel_def, eq:dom_im_def, eq:eq_class_def, eq:eq_class_rep}
\noeqref{eq:coord_order, eq:lex_order, eq:map_def, eq:comp_def}
\noeqref{eq:comp_assoc, eq:inv_map_def, eq:comp_inv}

\subsection{Бинарные отношения}

\begin{definition}
    Говорят, что между элементами множеств $X$ и $Y$ задано 
    \textdef{бинарное отношение} $R$, если фиксировано некоторое 
    подмножество их декартова произведения:
    \begin{equation} \label{eq:bin_rel_def}
        R \subseteq X \times Y \tag{1}
    \end{equation}
    В частном случае, когда $X = Y$, говорят, что отношение $R \subseteq X^2$ 
    задано \textdef{на множестве} $X$.
\end{definition}

\begin{designation}
    Принадлежность упорядоченной пары $(x, y)$ отношению $R$ 
    ($(x, y) \in R$) часто записывают в инфиксной форме: $x R y$. В 
    зависимости от контекста для бинарных отношений используются специальные 
    символы: $x \preceq y$, $x \sim y$, $x \prec y$ и другие.
\end{designation}

\begin{definition}
    Для любого бинарного отношения $R \subseteq X \times Y$ определяются:
    \begin{itemize}
        \item \textdef{Область определения} (domain):
        \item \textdef{Область значений} или \textdef{образ} (image):
    \end{itemize}
    \begin{equation} \label{eq:dom_im_def}
        \dom R = \{x \in X \mid \exists y \in Y : x R y\}, \qquad
        \im R = \{y \in Y \mid \exists x \in X : x R y\} \tag{2}
    \end{equation}
\end{definition}

\subsubsection{Свойства отношений на множестве}

\begin{definition}
    Пусть задано бинарное отношение $R \subseteq X^2$ на множестве $X$. 
    Определим фундаментальные свойства отношения $R$:
    \begin{enumerate}
        \item \textdef{Рефлексивность:} 
        $\forall x \in X : x R x$;
        \item \textdef{Иррефлексивность (антирефлексивность):} 
        $\forall x \in X : \neg(x R x)$;
        \item \textdef{Симметричность:} 
        $\forall x, y \in X : x R y \implies y R x$;
        \item \textdef{Антисимметричность:} 
        $\forall x, y \in X : (x R y \ \text{и} \ y R x) \implies x = y$ 
        (характерно для нестрогих порядков $\preceq$);
        \item \textdef{Асимметричность:} 
        $\forall x, y \in X : x R y \implies \neg(y R x)$ 
        (характерно для строгих порядков $\prec$);
        \item \textdef{Транзитивность:} 
        $\forall x, y, z \in X : (x R y \ \text{и} \ y R z) \implies x R z$.
    \end{enumerate}
\end{definition}

\subsection{Отношения эквивалентности}

\begin{definition}
    Бинарное отношение $R \subseteq X^2$ называется \textdef{отношением 
    эквивалентности} (обычно обозначается $\sim$), если оно одновременно 
    рефлексивно, симметрично и транзитивно.
\end{definition}

\begin{example} ~
    Приведём крайние примеры отношений эквивалентности на множестве $X$:
    \begin{enumerate}
        \item \emph{Полное (универсальное) отношение} $R = X^2$: $x R y$ 
        для любых $x, y \in X$ (все элементы эквивалентны друг другу).
        \item \emph{Отношение равенства (диагональ)} 
        $R = \{(x, x) \mid x \in X\}$: каждый элемент эквивалентен только 
        самому себе.
    \end{enumerate}
\end{example}

\begin{definition}
    \textdef{Классом эквивалентности} элемента $a \in X$ по отношению 
    эквивалентности $\sim$ называется подмножество всех элементов из $X$, 
    эквивалентных элементу $a$:
    \begin{equation} \label{eq:eq_class_def}
        \cl_\sim(a) = \{x \in X \mid x \sim a\} \tag{3}
    \end{equation}
\end{definition}

\begin{theorem}
    Класс эквивалентности порождается любым своим представителем:
    \begin{equation} \label{eq:eq_class_rep}
        b \in \cl_\sim(a) \implies \cl_\sim(b) = \cl_\sim(a) \tag{4}
    \end{equation}
\end{theorem}

\begin{proof}[Доказательство независимости класса от представителя]~\\
    \begin{quote}
        Пусть $b \in \cl_\sim(a)$, что по определению \eqref{eq:eq_class_def} 
        означает $b \sim a$. Докажем равенство множеств через взаимное 
        включение.

        1. Пусть $x \in \cl_\sim(b)$, тогда $x \sim b$. Так как $b \sim a$, 
        по свойству транзитивности получаем $x \sim a$, то есть 
        $x \in \cl_\sim(a)$. Следовательно, $\cl_\sim(b) \subseteq \cl_\sim(a)$.

        2. Пусть $x \in \cl_\sim(a)$, тогда $x \sim a$. Из условия $b \sim a$ 
        по свойству симметричности имеем $a \sim b$. Тогда из $x \sim a$ и 
        $a \sim b$ по транзитивности следует $x \sim b$, откуда 
        $x \in \cl_\sim(b)$. Значит, $\cl_\sim(a) \subseteq \cl_\sim(b)$.

        Из взаимного включения следует равенство $\cl_\sim(b) = \cl_\sim(a)$.
    \end{quote}
\end{proof}

\begin{theorem}
    Любые два класса эквивалентности либо не пересекаются, либо полностью 
    совпадают.
\end{theorem}

\begin{proof}[Доказательство о непересечении классов]~\\
    \begin{quote}
        Предположим, что два класса эквивалентности $\cl_\sim(a)$ и 
        $\cl_\sim(b)$ имеют хотя бы один общий элемент: 
        $x \in \cl_\sim(a) \cap \cl_\sim(b)$.

        По определению это означает, что $x \sim a$ и $x \sim b$. В силу 
        симметричности из $x \sim a$ следует $a \sim x$. Применяя 
        транзитивность к паре соотношений $a \sim x$ и $x \sim b$, получаем 
        $a \sim b$, то есть $a \in \cl_\sim(b)$. По теореме 
        \eqref{eq:eq_class_rep} отсюда немедленно следует, что 
        $\cl_\sim(a) = \cl_\sim(b)$.
    \end{quote}
\end{proof}

\begin{corollary}
    Всякое отношение эквивалентности разбивает исходное множество $X$ на 
    попарно непересекающиеся непустые классы эквивалентности (задаёт 
    \textdef{разбиение множества}).
\end{corollary}

\begin{definition}
    Множество всех классов эквивалентности множества $X$ по отношению 
    $\sim$ называется \textdef{фактор-множеством} и обозначается $X/{\sim}$ 
    (иногда также используется запись $\abs{X}_\sim$).
\end{definition}

\begin{remark}
    Переход от исходного множества элементов к фактор-множеству называется 
    \textdef{факторизацией}. Это важнейший математический приём, 
    позволяющий отождествлять схожие по изучаемому признаку объекты и 
    сводить бесконечные множества к конечным фактор-множествам (например, 
    арифметика остатков $\Z / m\Z$).
\end{remark}

\subsection{Отношения порядка}

\begin{definition}
    Бинарное отношение $\preceq$ на множестве $X$ называется 
    \textdef{отношением нестрогого частичного порядка}, если оно 
    рефлексивно, антисимметрично и транзитивно. Пара $\langle X, \preceq \rangle$ 
    называется \textdef{частично упорядоченным множеством (ЧУМ)}.
\end{definition}

\begin{example}
    Отношение $\le$ на числовых множествах $\N, \Z, \Q, \R$; отношение 
    включения $\subseteq$ на булеане $\mathcal{P}(X)$.
\end{example}

\begin{definition}
    Бинарное отношение $\prec$ на множестве $X$ называется 
    \textdef{отношением строгого частичного порядка}, если оно 
    иррефлексивно, асимметрично и транзитивно.
\end{definition}

\begin{example}
    Отношение $<$ для чисел; отношение строгого включения $\subset$ для 
    множеств.
\end{example}

\begin{remark}
    Термин «частичный» подчёркивает, что не любые два элемента множества $X$ 
    обязаны быть сравнимы между собой: могут существовать такие $x, y \in X$, 
    для которых не выполняется ни $x \preceq y$, ни $y \preceq x$. Если же 
    любые два элемента сравнимы ($\forall x, y \in X : x \preceq y$ или 
    $y \preceq x$), порядок называется \textdef{линейным} (или 
    \textdef{совершенным}).
\end{remark}

\begin{definition}[Способы упорядочивания кортежей]
    Само по себе выражение $(x_1, y_1) \preceq (x_2, y_2)$ на декартовом 
    произведении упорядоченных множеств требует уточнения способа сравнения. 
    Наиболее часто используются два канонических порядка:
    \begin{enumerate}
        \item \textdef{Покоординатный (паретовский) порядок:} сравнение 
        производится по каждой координате независимо:
        \begin{equation} \label{eq:coord_order}
            (x_1, y_1) \preceq (x_2, y_2) \iff 
            (x_1 \le x_2 \quad \text{и} \quad y_1 \le y_2) \tag{5}
        \end{equation}
        Строгая версия задаётся как $(x_1, y_1) \prec (x_2, y_2) \iff 
        (x_1, y_1) \preceq (x_2, y_2)$ и $(x_1, y_1) \neq (x_2, y_2)$. 
        Покоординатный порядок является частичным: например, пары $(1, 5)$ 
        и $(2, 3)$ несравнимы, так как $1 < 2$, но $5 > 3$.

        \item \textdef{Лексикографический (словарный) порядок:} сначала 
        сравниваются первые координаты, а при их равенстве --- вторые:
        \begin{equation} \label{eq:lex_order}
            (x_1, y_1) \preceq_{\mathrm{lex}} (x_2, y_2) \iff 
            (x_1 < x_2) \ \text{или} \ (x_1 = x_2 \ \text{и} \ y_1 \le y_2) \tag{6}
        \end{equation}
        Если исходные порядки на координатах линейны, то и 
        лексикографический порядок линеен: любые две пары сравнимы. 
        Например, $(1, 10) \prec_{\mathrm{lex}} (2, 0)$, так как $1 < 2$.
    \end{enumerate}
\end{definition}

\begin{definition}
    Пусть задано частично упорядоченное множество $\langle X, \preceq \rangle$.
    \begin{itemize}
        \item Элемент $x_{\min} \in X$ называется \textdef{минимальным}, 
        если в $X$ не существует элемента, строго меньшего его: 
        $\forall x \in X : (x \preceq x_{\min} \implies x = x_{\min})$. 
        Иными словами, любой элемент $x \in X$ либо несравним с $x_{\min}$, 
        либо удовлетворяет условию $x_{\min} \preceq x$.
        \item Элемент $x_{\text{наим}} \in X$ называется 
        \textdef{наименьшим}, если он сравним со всеми элементами множества 
        и меньше или равен каждому из них: 
        $\forall x \in X : x_{\text{наим}} \preceq x$.
    \end{itemize}
    Минимальных элементов в ЧУМ может быть несколько (или не быть вовсе). 
    Наименьший элемент, если он существует, всегда единственен и 
    одновременно является единственным минимальным элементом.
\end{definition}

\begin{definition}
    Частично упорядоченное множество $\langle X, \preceq \rangle$ называется 
    \textdef{фундированным} (или \textdef{нётеровым снизу}), если в любом 
    его непустом подмножестве существует хотя бы один минимальный элемент.
    Эквивалентное условие: в множестве $X$ не существует бесконечных 
    строго убывающих цепей вида $x_1 \succ x_2 \succ x_3 \succ \dots$.
\end{definition}

\begin{definition}
    Упорядоченное множество $\langle X, \preceq \rangle$ называется 
    \textdef{вполне упорядоченным}, если порядок $\preceq$ является 
    линейным и множество фундировано (то есть в любом его непустом 
    подмножестве существует наименьший элемент).
\end{definition}

\begin{example} ~
    \begin{itemize}
        \item Множество $\langle \N, \le \rangle$ фундировано и вполне 
        упорядочено: любое непустое подмножество натуральных чисел имеет 
        наименьший элемент, а любая убывающая цепь $5 > 4 > 3 > 2 > 1 > 0$ 
        обрывается за конечное число шагов.
        \item Множество $\langle \Z, \le \rangle$ не является фундированным: 
        подмножество отрицательных целых чисел не имеет минимального 
        элемента, и в нём существует бесконечная убывающая цепь 
        $0 > -1 > -2 > -3 > \dots$.
    \end{itemize}
\end{example}

\subsection{Отображения}

\begin{definition}
    Бинарное отношение $f \subseteq X \times Y$ называется 
    \textdef{отображением} (или \textdef{функцией}) из $X$ в $Y$, если 
    каждому элементу $x \in X$ соответствует не более одного элемента 
    $y \in Y$:
    \begin{equation} \label{eq:map_def}
        (x f y \quad \text{и} \quad x f z) \implies y = z \qquad 
        \forall x \in X, \ \forall y, z \in Y \tag{7}
    \end{equation}
    Обозначение: $f: X \to Y$, $y = f(x)$.
\end{definition}

\begin{definition}
    Отображение $f: X \to Y$ называется \textdef{полным (всюду 
    определённым)}, если $\dom f = X$, и \textdef{частичным}, если 
    $\dom f \subset X$. \emph{Далее по умолчанию рассматриваются только 
    полные отображения.}
\end{definition}

\subsubsection{Классификация отображений}

\begin{definition}
    Пусть задано полное отображение $f: X \to Y$. Оно называется:
    \begin{enumerate}
        \item \textdef{Инъективным} (\textdef{инъекцией}), если различные 
        элементы множества $X$ переходят в различные элементы множества $Y$:
        \[
            \forall x_1, x_2 \in X : 
            (x_1 \neq x_2 \implies f(x_1) \neq f(x_2)) \iff 
            (f(x_1) = f(x_2) \implies x_1 = x_2)
        \]
        \item \textdef{Сюръективным} (\textdef{сюръекцией}), если каждый 
        элемент множества $Y$ имеет хотя бы один прообраз в $X$ 
        ($\im f = Y$):
        \[
            \forall y \in Y \ \exists x \in X : f(x) = y
        \]
        \item \textdef{Биективным} (\textdef{биекцией} или \textdef{взаимно 
        однозначным отображением}), если оно одновременно инъективно и 
        сюръективно.
    \end{enumerate}
\end{definition}

\subsubsection{Композиция и обратные отображения}

\begin{definition}
    \textdef{Произведением (композицией)} отображений $g: X \to Y$ и 
    $f: Y \to Z$ называется отображение $f \circ g: X \to Z$, определяемое 
    правилом:
    \begin{equation} \label{eq:comp_def}
        (f \circ g)(x) = f(g(x)) \qquad \forall x \in X \tag{8}
    \end{equation}
\end{definition}

\begin{theorem}
    Композиция отображений ассоциативна. Для любых отображений 
    $g: X \to Y$, $f: Y \to Z$ и $h: Z \to W$ выполняется равенство:
    \begin{equation} \label{eq:comp_assoc}
        h \circ (f \circ g) = (h \circ f) \circ g \tag{9}
    \end{equation}
\end{theorem}

\begin{proof}[Доказательство ассоциативности композиции]~\\
    \begin{quote}
        Возьмём произвольный элемент $x \in X$. По определению композиции 
        \eqref{eq:comp_def} вычислим значение левой части равенства в точке $x$:
        \[
            \bigl(h \circ (f \circ g)\bigr)(x) = 
            h\bigl((f \circ g)(x)\bigr) = h\bigl(f(g(x))\bigr)
        \]
        Аналогично вычислим значение правой части в точке $x$:
        \[
            \bigl((h \circ f) \circ g\bigr)(x) = 
            (h \circ f)\bigl(g(x)\bigr) = h\bigl(f(g(x))\bigr)
        \]
        Оба отображения совпадают в каждой точке $x \in X$, следовательно, 
        они равны.
    \end{quote}
\end{proof}

\begin{definition}
    Для отображения $f: X \to Y$ отображение $f^{-1}: Y \to X$ называется 
    \textdef{обратным}, если одновременно выполняются соотношения:
    \begin{equation} \label{eq:inv_map_def}
        \begin{cases}
            f \circ f^{-1} = e_Y, \\
            f^{-1} \circ f = e_X,
        \end{cases} \tag{10}
    \end{equation}
    где $e_X: X \to X$ и $e_Y: Y \to Y$ --- тождественные отображения на 
    множествах $X$ и $Y$ соответственно ($e_X(x) = x$, $e_Y(y) = y$). 
    Отображение, имеющее обратное, называется \textdef{обратимым}.
\end{definition}

\begin{remark}
    Далеко не каждое отображение имеет обратное. Условие существования 
    обратного отображения устанавливается следующими утверждениями.
\end{remark}

\begin{lemma} \label{lem:inv_inj_surj}
    Пусть даны отображения $f: X \to Y$ и $g: Y \to X$, причём 
    $g \circ f = e_X$. Тогда отображение $f$ инъективно, а отображение $g$ 
    сюръективно.
\end{lemma}

\begin{proof}[Доказательство леммы]~\\
    \begin{quote}
        1. \textbf{Инъективность $f$.} Докажем от противного. Пусть 
        существуют $x_1, x_2 \in X$ ($x_1 \neq x_2$), такие что 
        $f(x_1) = f(x_2)$. Применим к обеим частям равенства отображение $g$:
        \[
            g(f(x_1)) = g(f(x_2)) \implies 
            (g \circ f)(x_1) = (g \circ f)(x_2)
        \]
        По условию $g \circ f = e_X$, откуда $e_X(x_1) = e_X(x_2)$, то есть 
        $x_1 = x_2$. Получено противоречие, значит, $f$ инъективно.

        2. \textbf{Сюръективность $g$.} Покажем, что 
        $\forall x \in X \ \exists y \in Y : g(y) = x$. Для произвольного 
        $x \in X$ положим $y = f(x) \in Y$. Тогда:
        \[
            g(y) = g(f(x)) = (g \circ f)(x) = e_X(x) = x
        \]
        Искомый прообраз $y$ найден, следовательно, $g$ сюръективно.
    \end{quote}
\end{proof}

\begin{theorem}[\textbf{Критерий обратимости отображения}]
    Отображение $f: X \to Y$ обратимо тогда и только тогда, когда оно 
    биективно.
\end{theorem}

\begin{proof}[Доказательство критерия обратимости]~\\
    \begin{quote}
        Докажем теорему в обе стороны.

        $\implies$ \textbf{Необходимость (обратимость влечёт биективность).}\\
        Пусть отображение $f$ обратимо, тогда существует $f^{-1}: Y \to X$, 
        удовлетворяющее системе \eqref{eq:inv_map_def}. Из равенства 
        $f^{-1} \circ f = e_X$ по лемме \ref{lem:inv_inj_surj} следует 
        инъективность $f$. Из второго равенства $f \circ f^{-1} = e_Y$ по 
        той же лемме \ref{lem:inv_inj_surj} следует сюръективность $f$. 
        Следовательно, $f$ биективно.

        $\impliedby$ \textbf{Достаточность (биективность влечёт обратимость).}\\
        Пусть $f$ биективно. В силу сюръективности для каждого $y \in Y$ 
        существует прообраз $x \in X$, такой что $f(x) = y$, а в силу 
        инъективности такой элемент $x$ единственен. Положим по определению 
        $f^{-1}(y) = x$. Проверим выполнение условий \eqref{eq:inv_map_def}:
        \begin{align*}
            (f \circ f^{-1})(y) &= f(f^{-1}(y)) = f(x) = y = e_Y(y) 
            \implies f \circ f^{-1} = e_Y, \\
            (f^{-1} \circ f)(x) &= f^{-1}(f(x)) = f^{-1}(y) = x = e_X(x) 
            \implies f^{-1} \circ f = e_X.
        \end{align*}
        Таким образом, отображение $f^{-1}$ корректно определено и является 
        обратным к $f$.
    \end{quote}
\end{proof}

\begin{theorem}[\textbf{Свойства композиции отображений}] ~
    \begin{enumerate}
        \item Композиция инъективных отображений инъективна.
        \item Композиция сюръективных отображений сюръективна.
        \item Композиция биективных отображений биективна.
    \end{enumerate}
\end{theorem}

\begin{proof}[Доказательство свойств композиции]~\\
    \begin{quote}
        Пусть даны отображения $g: X \to Y$ и $f: Y \to Z$. Рассмотрим их 
        композицию $f \circ g: X \to Z$.

        1. Пусть $g$ и $f$ инъективны, и пусть $(f \circ g)(x_1) = 
        (f \circ g)(x_2)$, то есть $f(g(x_1)) = f(g(x_2))$. В силу 
        инъективности $f$ отсюда следует $g(x_1) = g(x_2)$, а в силу 
        инъективности $g$ получаем $x_1 = x_2$. Значит, $f \circ g$ 
        инъективна.

        2. Пусть $g$ и $f$ сюръективны. Возьмём произвольный элемент 
        $z \in Z$. Так как $f$ сюръективна, найдётся $y \in Y$, такой что 
        $f(y) = z$. Так как $g$ сюръективна, для найденного $y$ существует 
        $x \in X$, такой что $g(x) = y$. Тогда $(f \circ g)(x) = f(g(x)) = 
        f(y) = z$, что доказывает сюръективность $f \circ g$.

        3. Биективность композиции биекций непосредственно следует из 
        пунктов 1 и 2.
    \end{quote}
\end{proof}

\begin{theorem}
    Если отображения $g: X \to Y$ и $f: Y \to Z$ обратимы, то их 
    композиция $f \circ g: X \to Z$ также обратима, причём:
    \begin{equation} \label{eq:comp_inv}
        (f \circ g)^{-1} = g^{-1} \circ f^{-1} \tag{11}
    \end{equation}
\end{theorem}

\begin{proof}[Доказательство формулы обратного к композиции]~\\
    \begin{quote}
        Обозначим $h = g^{-1} \circ f^{-1}: Z \to X$. Используя 
        ассоциативность композиции \eqref{eq:comp_assoc}, проверим условия 
        обратного отображения \eqref{eq:inv_map_def}:
        \begin{align*}
            (f \circ g) \circ h &= (f \circ g) \circ (g^{-1} \circ f^{-1}) = 
            f \circ (g \circ g^{-1}) \circ f^{-1} = 
            f \circ e_Y \circ f^{-1} = f \circ f^{-1} = e_Z, \\
            h \circ (f \circ g) &= (g^{-1} \circ f^{-1}) \circ (f \circ g) = 
            g^{-1} \circ (f^{-1} \circ f) \circ g = 
            g^{-1} \circ e_Y \circ g = g^{-1} \circ g = e_X.
        \end{align*}
        Следовательно, $h = g^{-1} \circ f^{-1}$ является обратным 
        отображением к $f \circ g$.
    \end{quote}
\end{proof}

\begin{theorem}
    Обратное отображение (если оно существует) единственно.
\end{theorem}

\begin{proof}[Доказательство единственности обратного отображения]~\\
    \begin{quote}
        Пусть у отображения $f: X \to Y$ существуют два обратных 
        отображения $g_1: Y \to X$ и $g_2: Y \to X$. Тогда по определению 
        \eqref{eq:inv_map_def} выполняются равенства $g_1 \circ f = e_X$ и 
        $f \circ g_2 = e_Y$. Используя ассоциативность композиции, запишем 
        цепочку тождеств:
        \[
            g_1 = g_1 \circ e_Y = g_1 \circ (f \circ g_2) = 
            (g_1 \circ f) \circ g_2 = e_X \circ g_2 = g_2
        \]
        Таким образом, $g_1 = g_2$.
    \end{quote}
\end{proof}

\subsubsection{Алгебраические структуры отображений}

Тождественное отображение играет роль нейтрального элемента: \(f \circ e_X = e_Y \circ f = f\).

Если рассматривать множество всех отображений из множества в само себя (\(f: X \to X\)), мы получаем следующие алгебраические структуры:
\begin{enumerate}
    \item \textdef{Полугруппа} — множество отображений с ассоциативной операцией композиции.
    \item \textdef{Моноид} — полугруппа, в которой есть нейтральный элемент (\(e_X\)).
    \item \textdef{Группа} — если рассматривать \emph{только биективные} отображения \(X \to X\), то каждое из них имеет обратное (\(f^{-1}\)). Моноид с обратимыми элементами образует группу.
\end{enumerate}